%!TEX root = forallxyyc.tex
\part{逻辑的核心概念}
\label{ch.intro}
\addtocontents{toc}{\protect\mbox{}\protect\hrulefill\par}

\chapter{论证}
\label{s:Arguments}

正如在~\hyperref[preface]{前言}中提到的，逻辑的用途\emph{很多}。我们在这里重点关注如何用它来评估论证：区分好的和坏的论证。

在日常语言中，我们有时用“论证”这个词来指激烈的争吵。如果你和一位朋友发生了这种意义上的“论证”，那说明你们之间情况不妙。逻辑不关心这种咬牙切齿、互相拉扯的场面。它们并非我们意义上的论证；它们仅仅是分歧。

按照我们的理解，论证更像是下面这种情形：

\begin{earg}\label{argButlerGardner}
		\item 要么是管家干的，要么是园丁干的。
		\item 不是管家干的。
		\item[\texttherefore] 所以是园丁干的。
	\end{earg}

这里我们有一个句子序列。论证第三行上的三个点读作“所以”。它们表明最后一个句子表达的是论证的\emph{结论}。在此之前的两个句子是论证的\emph{前提}。如果你相信这些前提，并且认为结论是从这些前提得出的——即，这个论证是有效的——那么这（或许）就给了你一个相信该结论的理由。

逻辑学家感兴趣的正是这类东西。我们将论证定义为：一系列前提与一个结论的组合。

本部分讨论一些适用于像英语这类自然语言中论证的基本逻辑概念。一开始就清晰地理解什么是论证以及论证有效意味着什么，这一点很重要。后续我们将用一种形式语言来刻画自然语言中的论证。%We want formal validity, as defined in the formal language, to have at least some of the important features of natural-language validity.

在刚举的例子中，我们用独立的句子分别表达了论证的两个前提，并用第三个句子表达了结论。许多论证都以这种方式表达，但一个单独的句子也可以包含一个完整的论证。例如：

\begin{quote}
		 管家有不在场证明；所以他不可能是凶手。
	\end{quote}

这个论证包含一个前提，后跟一个结论。

许多论证以前提开头，以结论结尾，但并非总是如此。本节开头的那个论证，同样可以把结论放在开头来呈现：

\begin{quote}
		是园丁干的。毕竟，凶手要么是管家，要么是园丁。而且管家没干。
	\end{quote}

同样，它也可能把结论放在中间：

\begin{quote}
		不是管家干的。因此，既然是园丁要么是管家干的，所以是园丁干的。
	\end{quote}

面对一个论证时，我们想知道结论是否从前提得出。所以，首先要做的是区分出前提和结论。作为指引，以下词语常用来指示论证的结论：

\begin{center}
		所以，因此，因而，于是，相应地，从而
	\end{center}

因此，它们有时被称为\define{结论提示词（conclusion indicator words）}。

相比之下，以下表达式是\define{前提提示词（premise indicator words）}，因为它们通常指示其后的内容为前提，而非结论：

\begin{center}
		既然，因为，鉴于
	\end{center}

但在分析论证时，良好的判断力无可替代。

\newglossaryentry{premise indicator word}
{
name=前提指示语,
description={用来指示接下来是论证前提的词语或短语，如“因为”}
}

\newglossaryentry{conclusion indicator word}
{
name=结论指示语,
description={用来指示接下来是论证结论的词语或短语，如“所以”}
}

\newglossaryentry{argument}
{
name=论证,
description={分为\gls{premise}和\gls{conclusion}的一系列相互关联的句子}
}

\newglossaryentry{premise}
{
name=前提,
description={论证中除\gls{conclusion}外的句子}
}

\newglossaryentry{conclusion}
{
name=结论,
description={论证的最后一个句子}
}

\section{句子}
\label{intro.sentences}

为了追求完全一般性，我们可以将\define{论证（argument）}定义为一个句子序列。序列开始的句子是前提，序列的最后一个句子是结论。如果前提为真并且论证是好的，那么你就有理由接受该结论。

在逻辑学中，我们只对那些能够充当论证前提或结论的句子感兴趣，即那些可以为真或为假的句子。因此，我们将自己限于这类句子，并将一个\define{句子（sentence）}定义为可以为真或为假的句子。

不应将“句子可以为真或为假”这一观念与“事实和观点”之间的区别相混淆。逻辑中的句子通常表达那些可算作事实的事情——例如“鲁道夫·卡尔纳普生于龙斯多夫”或“西蒙娜·德·波伏瓦喜欢散步”。它们也可以表达那些你可能认为是观点的事情——例如，“大黄很好吃”。换言之，一个句子不会因为我们不知道其真假，或因为其真假是观点问题，而被取消作为论证一部分的资格。只要它是那种可以为真或为假的句子，它就能扮演前提或结论的角色。

此外，有一些在语言学或语法课中算作“句子”的东西，在我们这里不会被算作逻辑中的句子。

\paragraph{疑问句} 在语法课上，“你困了吗？”会被算作一个疑问句。虽然你可能困了或者很清醒，但这个问题本身既不真也不假。因此，疑问句在逻辑中不被算作句子。假设你回答这个问题：“我不困。”这个回答要么真要么假，所以它在逻辑意义下是一个句子。一般而言，\emph{疑问句}不算句子，但\emph{回答}算。

“这门课是关于什么的？”不是（我们意义上的）句子。“没人知道这门课是关于什么的”是一个句子。

\paragraph{祈使句} 命令常以祈使句形式表达，如“醒来！”、“坐直”等。在语法课上，这些会被算作祈使句。尽管坐直对你来说可能是好事也可能不是，但这个命令本身既不真也不假。但要注意，命令并不总是以祈使句形式表达。“你会尊重我的权威”\emph{是}要么真要么假的——你要么会，要么不会——因此它在逻辑意义下算作一个句子。

\paragraph{感叹句} “哎哟！”有时被称为感叹句，但它既不真也不假。我们会把“哎哟，我伤到脚趾了！”看作和“我伤到脚趾了”意思相同。“哎哟”并没有增添任何可以为真或为假的内容。

\practiceproblems
在一些章节的末尾，会有回顾和探索本章内容的练习题。实际动手做题是无法替代的，因为学习逻辑更多地是培养一种思维方式，而非记忆事实。

\medskip

下面是第一个练习。标出下列每个论证中表达结论的短语：

\begin{compactlist}
	\item 天晴了。所以我应该戴上太阳镜。
	\item 当时一定是个晴天。毕竟我确实戴了太阳镜。
	\item 除了你，没人碰过饼干罐。而犯罪现场撒满了饼干屑。你就是罪犯！
	\item 案发时，斯卡利特小姐和普拉姆教授在书房。
	格林牧师拿着烛台在舞厅，而且我们知道他手上没有血迹。因此，
	芥末上校是在厨房用铅管干的。别忘了，毕竟枪没有开火。
\end{compactlist}

\chapter{逻辑的范围}
\label{s:Valid}

\section{后承与有效性}

在 \cref{s:Arguments} 中，我们讨论了论证，即由一组句子（前提）后跟一个单独的句子（结论）构成的东西。我们说过，像“所以”这样的一些词语指示哪个句子是结论。“所以”自然暗示了前提和结论之间有一种联系，即结论 \emph{从前提得出}，或者说 \emph{是前提的后承}。

后承这一概念是逻辑学关注的核心之一。甚至可以说，逻辑学是关于“什么从什么得出”的科学。逻辑学发展出的理论和工具，能告诉我们一个句子何时从另一些句子得出。

那么，\cref{s:Arguments} 中讨论的主要论证又如何呢？

\begin{earg}
	\item[] 要么是管家干的，要么是园丁干的。
	\item[] 不是管家干的。
	\item[\texttherefore] 是园丁干的。
\end{earg}

我们不知道这个论证中的句子具体指什么。也许你怀疑这里的“做了”指的是某桩未具体说明的罪行。你可能想象这个论证出自一部推理小说或电视剧，或许是一位正在梳理证据的侦探所说。但即便没有任何这类背景信息，你多半也会认为这是一个好论证：因为无论这些前提具体指什么，只要它们都为真，其结论就不可能为假。如果第一个前提为真，即“管家做了或园丁做了”为真，那么至少他们中有一人“做了”，无论那意味着什么。而如果第二个前提为真，即管家没有“做”，那么剩下的就只有一种可能：“园丁做了”必然为真。此时，结论是从前提得出的。我们称具有这种性质的论证为 \define{有效的}。

作为对比，考虑以下论证：

\begin{earg}\label{argMaidDriver}
	\item[] 如果司机做了，那么女佣没有做。
	\item[] 女佣没有做。
	\item[\texttherefore] 司机做了。
\end{earg}

我们依然不知道这里谈论的是什么。但是，你很可能也同意，这个论证与上一个在一个重要方面存在区别：即使其前提为真，也\emph{不能保证}其结论为真。这个论证的前提本身，并未排除司机或女佣之外的人“做了”的可能性。因此，存在一种情形：两个前提都为真，然而司机并未作案——即结论不为真。在这第二个论证中，结论并非由前提得出。如果一个论证像此例一样，其结论并非由前提得出，我们就称它是 \define{无效的}。

\section{情形与有效性的类型}

我们是如何判定第二个论证无效的呢？我们指出了一个前提为真而结论不为真的情形。该情形就是：司机和女佣都没有作案，而是某个第三人做了。我们称这样的情形为该论证的一个 \define{反例}。如果一个论证存在反例，那么其结论就不可能是前提的后果。因为，结论要成为前提的后果，前提的真必须能确保结论的真。如果存在反例，前提的真就无法确保结论的真。

作为逻辑学家，我们希望有能力判定一个论证的结论何时能从前提得出。并且，当没有反例——即不存在所有前提都为真而结论为假的情形时，结论就是前提的后果。这引出了如下定义：

	\factoidbox{
		一个语句 $A$ 是语句 $B_1$, \dots, $B_n$ 的 \define{后果}，当且仅当不存在 $B_1$, \dots, $B_n$ 都真而 $A$ 为假的情形。（此时我们也说 $A$ 从 $B_1$, \dots, $B_n$ \define{得出}，或者说 $B_1$, \dots, $B_n$ \define{蕴涵}~$A$。）
	}

这个“定义”并不完整：它没有告诉我们“情形”是什么，也没有说明“在一个情形中为真”意味着什么。截至目前，我们只看到了一个例子：一个涉及三名角色的设想场景。在该场景中——包括司机、女佣和某个第三人——司机和女佣没有作案，但第三人做了。如所描述，在这一场景中，司机没有作案，因此它是“司机做了”这一语句为假的一个情形。我们第二个论证的前提为真，但结论为假：该场景构成一个反例。

我们说过，结论是前提后果的论证称为有效的，结论不是前提后果的论证称为无效的。既然我们现在至少初步定义了“后果”，我们将其记录如下：

	\factoidbox{


\begin{factoiditems}
			\item 一个论证是 \define{有效的} 当且仅当其结论是前提的后果。
			\item 一个论证是 \define{无效的} 当且仅当它不是有效的，即它存在一个反例。
		\end{factoiditems}

}

\newglossaryentry{valid}
{
name=有效,
description={论证具有的一种属性，其结论是前提的后果}
}

\newglossaryentry{invalid}
{
name=无效,
description={当结论不是前提的后果时，论证所具有的一种属性；与 \gls{valid} 相对}
}

逻辑学的任务之一，就是使“情形”这一概念更加精确，并探究当“情形”以特定方式被精确化时，哪些论证是有效的。如果我们把“情形”理解为（诸如第二个论证的反例那样的）“设想情形”，那么显然第一个论证是有效的。如果我们设想一种情形：或者管家或者园丁做了那件事，并且管家没有做，那么我们就自动地设想了一种园丁做了那件事的情形。因此，任何使我们第一个论证的前提为真的设想情形，都自动地使其结论为真。这使得第一个论证成为有效的。

将“情形”具体解释为“设想情形”是一种推进，但这并非终点。第一个问题在于，我们不知道什么才算作一个设想情形。它们受物理定律限制吗？还是受制于最广泛意义上的、可设想的东西？我们对这些问题的回答，决定了哪些论证会被我们算作有效的。

假设对第一个问题的回答是“是”。考虑以下论证：


\begin{earg}
		\item[] 飞船“罗西南迪号”从泰科空间站到木星花了六个小时。
		\item[\texttherefore] 泰科空间站与木星之间的距离小于 140 亿公里。
	\end{earg}


这个论证的一个反例会是这样一种情形：“罗西南迪号”在六小时内完成了超过 140 亿公里的航程，即超过了光速。由于这种情形与物理定律不相容，如果设想情形必须符合物理定律，那么就不存在这样的情形。然而，如果设想情形不受物理定律限制，那么反例就是存在的：“罗西南迪号”超光速飞行的情形。

假设对第二个问题的回答是“是”，再考虑另一个论证：


\begin{earg}
		\item[] Priya 是一名 ophthalmologist。
		\item[\texttherefore] Priya 是一名 eye doctor。
	\end{earg}


如果我们只允许可设想的情形，那么这个论证也是有效的。如果你设想 Priya 是一名 ophthalmologist，你也就由此设想了 Priya 是一名 eye doctor。这正是“ophthalmologist”和“eye doctor”这两个词的含义。一个 Priya 是 ophthalmologist 但不是 eye doctor 的情形，会被这两个词之间的概念联系排除掉。

根据我们将哪些类型的情形视为潜在的反例，我们会得到不同的后果与有效性概念。有效性是以没有反例为特征的。但对于不同的有效性概念，什么才算作反例可能有所不同。例如，我们可能会排除违反自然法则的反例，或者排除违反像“eye doctor”和“ophthalmologist”这类词汇之间概念联系的反例。因此，如果一个论证没有违反自然法则的反例，我们可以称它为 \define{法则有效的}；如果一个论证没有违反词汇间概念联系的反例，则可以称它为 \define{概念有效的}。对于这两种有效性概念而言，世界的某些方面（例如，自然法则是什么）以及论证中语句意义的某些方面（例如，“ophthalmologist”就是指某种 eye doctor），都会影响该论证是否被视为有效。

\section{形式有效性}\label{s:FormalValidity}

然而，\emph{逻辑}后果的一个显著特征是，它不应依赖于前提和结论的内容，而只应依赖于它们的逻辑形式。换句话说，作为逻辑学家，我们希望发展一种理论，来捕捉论证成为有效的另一种方式。例如，下面两个论证：

\begin{earg}
	\item[] Priya 要么是 ophthalmologist，要么是 dentist。
	\item[] Priya 不是 dentist。
	\item[\texttherefore] Priya 是 eye doctor。
\end{earg}

和

\begin{earg}
	\item[] Priya 要么是 ophthalmologist，要么是 dentist。
	\item[] Priya 不是 dentist。
	\item[\texttherefore] Priya 是 ophthalmologist。
\end{earg}

都是有效论证。但第一个论证的有效性依赖于内容（即“ophthalmologist”和“eye doctor”的意义），而第二个则不依赖。第二个论证是 \define{形式上有效的}。我们可以将这个论证的“形式”描述为一种模式，大致如下：

\begin{earg}
    \item[] $A$ 要么是一个 $X$，要么是一个 $Y$。
    \item[] $A$ 不是一个 $Y$。
    \item[\texttherefore] $A$ 是一个 $X$。
\end{earg}

此处的 $A$、$X$ 和 $Y$ 是占位符，代表合适的表达式；将其代入对应位置后，就能将这一模式转变成一个由句子构成的论证。例如，

\begin{earg}
    \item[] 梅要么是数学家，要么是植物学家。
    \item[] 梅不是植物学家。
    \item[\texttherefore] 梅是数学家。
\end{earg}

就是一个具有相同形式的论证，但本节开头给出的第一个论证却不属于这种形式：我们需要将 $Y$ 替换为不同的表达式（一次是“ophthalmologist”，一次是“eye doctor”），才能从该模式得到它。

此外，第一个论证并不是形式有效的。\emph{它的}形式是这样的：

\begin{earg}
    \item[] $A$ 要么是一个 $X$，要么是一个 $Y$。
    \item[] $A$ 不是一个 $Y$。
    \item[\texttherefore] $A$ 是一个 $Z$。
\end{earg}

在这个模式中，我们可以用“ophthalmologist”替换 $X$，用“eye doctor”替换 $Z$，从而得到原论证。但下面这个论证也具有相同形式：

\begin{earg}
    \item[] 梅要么是数学家，要么是植物学家。
    \item[] 梅不是植物学家。
    \item[\texttherefore] 梅是杂技演员。
\end{earg}

这个论证显然是无效的，因为完全可以想象一位名叫梅的数学家并非杂技演员。

作为逻辑学家，我们的策略是提出一种“情形”概念，使得一个论证仅当它是形式有效时，才会在这种“情形”概念下被判定为有效。显然，这种“情形”概念将不仅违反某些自然法则，还会违反英语的某些语义法则。由于第一个论证在这种意义下是无效的，我们必须允许一种构成反例的情形：Priya 是一位 ophthalmologist 但不是一位 eye doctor。这种情形是不可设想的：它被“ophthalmologist”和“eye doctor”的含义所排除。

当我们为了评估论证的有效性而考虑各种情形时，会做一个重要假设：我们假设每一种情形都使所考虑的每一个句子为真或不真。这首先意味着，任何无法确定我们论证中某句子真假的想象场景，都不会被视为潜在的反例。例如，一个 Priya 是牙医但不是眼科医生的场景，在考虑本节前几个论证时会被视为一种可考虑的情形，但在考虑最后两个论证时则不会：它没有告诉我们梅是否是数学家、植物学家或杂技演员。如果一个情形没有使一个句子为真，我们就说它使这个句子为\define{假}。因此，我们假设情形使句子为真或为假，但不会同时既为真又为假。\footnote{句子非真即假且不兼而有之的假设，称为“二值性”。即便此假设在你看来是常识，它在逻辑哲学中仍有争议。首先，有些逻辑学家希望考虑句子既不为真也不为假，而是具有某种中间真值的情形。更具争议的是，一些哲学家认为应允许句子同时既真又假的可能性。存在允许句子既不为真也不为假、或同时既真又假的逻辑系统，但本书不予讨论。}

\section{可靠的论证}

在继续执行这一策略前，需稍作澄清。我们将论证定义为一个句子集合，其中一句（结论）应当从其他句子（前提）得出。这种意义上的论证在日常和科学论述中一直使用。此时，论证被用来支持甚至证明其结论。如果一个论证是有效的，它将支持其结论，但\emph{仅当}其所有前提都为真。有效性排除了前提为真而结论同时不为真的可能性。它本身并不排除结论不为真的可能性。换言之，一个有效论证完全可能拥有一个不真的结论！

考虑下例：

\begin{earg}
    \item[] 橘子要么是水果，要么是乐器。
    \item[] 橘子不是水果。
    \item[\texttherefore] 橘子是乐器。
\end{earg}

此论证的结论是荒谬的。然而，它确实是从前提得出的。\emph{如果}两个前提都为真，\emph{那么}结论必然为真。所以该论证是有效的。

反之，前提为真且结论为真，并不足以使论证有效。考虑下例：

\begin{earg}
    \item[] 伦敦在英格兰。
    \item[] 北京在中国。
    \item[\texttherefore] 巴黎在法国。
\end{earg}

此论证的前提与结论事实上均为真，但该论证是无效的。倘若巴黎宣布从法国独立，则结论将不再为真，尽管两个前提仍保持为真。因此，存在一种情形，该论证的前提为真而结论不为真。所以，此论证无效。

需牢记的要点是：有效性无关论证中句子的\emph{实际}真假。它关乎是否\emph{可能}出现所有前提为真而结论同时不为真的情形（于某个假想情形中）。\emph{事实上}如何并不起特殊作用；事实本身不能决定论证是否有效。\footnote{诚然，有一种例外：若前提事实上为真且结论事实上不为真，则我们正身处一个反例中；故此论证无效。}事物实际状况的任何方面，其本身都不能决定论证是否有效。常言道逻辑不关心感受。实则，它也不关心事实。

那么，当我们用论证来证明其结论\emph{为真}时，需要两样东西：首先，论证需有效；即结论必须能从前提得出。其次，前提本身需为真。我们将称一个论证是\define{可靠（sound）}的，当且仅当它既有效且所有前提均为真。

\newglossaryentry{sound}
{
name=可靠,
description={论证的一个属性，成立的条件是该论证有效且所有前提为真}
}

反过来，当你想反驳一个论证时，有两种选择：指出（一个或多个）前提不真，或指出论证无效。然而，逻辑仅能助你处理后一种情况！

\section{归纳论证}

许多好的论证是无效的。考虑以下论证：

\begin{earg}
    \item[] 迄今为止的每一个冬天，卡尔加里都下过雪。
    \item[\texttherefore] 即将到来的这个冬天，卡尔加里也会下雪。
\end{earg}

这个论证从对许多（过去）情形的观察概括出关于所有（未来）情形的结论。此类论证称为\define{归纳（inductive）}论证。然而，它是无效的。即便卡尔加里迄今每年冬天都下雪，那里在整个即将到来的冬天保持干燥仍是\emph{可能}的。事实上，即使卡尔加里从此每年冬天都下雪，我们仍可\emph{想象}一种情形：今年是首个全冬无雪的年份。此假想情形正是该论证前提为真而结论不为真的一个案例，从而使该论证无效。

关键在于，归纳论证——即便是好的归纳论证——也非（演绎）有效。它们并非\emph{滴水不漏}。尽管可能性不大，但即使所有前提都为真，其结论也\emph{可能}为假。本书将（完全）搁置“何谓一个好的归纳论证”之问题。我们的兴趣仅在区分（演绎）有效的论证与无效的论证。

因此：我们关注一个结论是否从某些前提\emph{得出}。但请勿说前提\emph{推断出}结论。蕴涵是前提与结论间的一种关系；推理则是我们之所为。故当结论能从前提得出时，若欲提及推理，可以说\emph{人们可从前提推断出}该结论。

\practiceproblems
\problempart
下列论证中哪些有效？哪些无效？

\begin{compactlist}
\item
\begin{earg}
\item 苏格拉底是人。
\item 所有人都是胡萝卜。
\item[\texttherefore] 苏格拉底是胡萝卜。
\end{earg}
\item
\begin{earg}
\item Abe Lincoln 要么出生在伊利诺伊州，要么他曾是总统。
\item Abe Lincoln 从未当过总统。
\item[\texttherefore] Abe Lincoln 出生在伊利诺伊州。
\end{earg}
\item
\begin{earg}
\item 如果我扣动扳机，Abe Lincoln 就会死。
\item 我没有扣动扳机。
\item[\texttherefore] Abe Lincoln 不会死。
\end{earg}
\item
\begin{earg}
\item Abe Lincoln 要么来自法国，要么来自卢森堡。
\item Abe Lincoln 不是来自卢森堡。
\item[\texttherefore] Abe Lincoln 来自法国。
\end{earg}
\item
\begin{earg}
\item 如果今天世界末日，那么我明天早上就不需要起床。
\item 我明天早上需要起床。
\item[\texttherefore] 今天不会是世界末日。
\end{earg}
\item
\begin{earg}
\item Joe 现在 19 岁。
\item Joe 现在 87 岁。
\item[\texttherefore] Bob 现在 20 岁。
\end{earg}
\end{compactlist}

\problempart
\label{pr.EnglishCombinations}
是否存在：

\begin{compactlist}
		\item 一个有效论证，它有一个假前提和一个真前提？
		\item 一个有效论证，它的所有前提都是假的？
		\item 一个有效论证，它的所有前提和结论都是假的？
		\item 一个无效论证，可以通过增加一个新前提而变成有效的？
		\item 一个有效论证，可以通过增加一个新前提而变成无效的？
	\end{compactlist}

在每种情况下：如果存在，请举例说明；如果不存在，请解释原因。

\chapter{其他逻辑概念}\label{s:BasicNotions}

在 \cref{s:Valid} 中，我们引入了后承的概念，即论证的有效性。这是逻辑学中最重要的概念之一。在本节中，我们将引入一些同样重要的概念。这些概念，和有效性一样，都依赖于“语句在情形中为真（或不为真）”这一思想。在本节的后续部分，我们将把“情形”理解为可设想的场景，也就是我们用来定义概念有效性时使用的那个意义。关于不同的有效性类型我们曾做出的区分，也可以类似地适用于我们引入的新概念：如果我们使用不同的“情形”概念，就会得到不同的概念。并且，作为逻辑学家，我们最终会考虑一种比这里所采用的更宽泛的“情形”定义。

%\section{Truth values}
%As we said in \cref{s:Arguments}, arguments consist of premises and a conclusion. Note that many kinds of English sentence cannot be used to express premises or conclusions of arguments. For example:
%	\begin{itemize}
%		\item \textbf{Questions}, e.g.\ “are you feeling sleepy?'
%		\item \textbf{Imperatives}, e.g.\ “Wake up!'
%		\item \textbf{Exclamations}, e.g.\ “Ouch!'
%	\end{itemize}
%The common feature of these three kinds of sentence is that they are not \emph{assertoric}: they cannot be true or false. It does not even make sense to ask whether a \emph{question} is true (it only makes sense to ask whether the \emph{answer} to a question is true).

%The general point is that, the premises and conclusion of an argument must be capable of having a \define{truth value}. The two truth values that concern us are just True and False.

\section{共同可满足性}

考虑下面两个句子：

\begin{compactlist}
		\item[B1.] Jane 唯一的兄弟比她矮。
		\item[B2.] Jane 唯一的兄弟比她高。
	\end{compactlist}

单凭逻辑无法告诉我们这两个句子中哪个为真（如果其中确有真句的话）。然而我们可以说，\emph{如果}第一个句子（B1）为真，\emph{那么}第二个句子（B2）必定为假。类似地，如果 B2 为真，那么 B1 必定为假。不存在一种可能的情形能让这两个句子同时为真。这两个句子彼此不相容，它们不可能同时为真。这引出了以下定义：
	\factoidbox{
		句子是 \define{共同可满足的（jointly possible）}，当且仅当存在一个情形使得它们全部为真。
	}
B1 和 B2 是 \emph{共同不可满足的}；而下面两个句子，例如，就是共同可满足的：

\begin{compactlist}
		\item[B1.] Jane 唯一的兄弟比她矮。
		\item[B2.] Jane 唯一的兄弟比她年轻。
	\end{compactlist}

\newglossaryentry{possibility}
{
name=共同可满足性,
text={共同可满足的},
description={一些句子在单个情形下全部为真时所具有的性质}
}

我们可以询问任意数量句子的共同可满足性。例如，考虑以下四个句子：

\begin{compactlist}
		\item[G1.] \label{MartianGiraffes} 野生动物园里至少有四只长颈鹿。
		\item[G2.] 野生动物园里恰好有七只大猩猩。
		\item[G3.] 野生动物园里的火星人不超过两个。
		\item[G4.] 野生动物园里的每一只长颈鹿都是火星人。
	\end{compactlist}

G1 和 G4 共同蕴涵了公园里至少有四只火星长颈鹿。这与 G3 矛盾，因为 G3 意味着那里最多只有两只火星长颈鹿。所以句子 G1--G4 是共同不可满足的。它们不可能全部为真。（注意，句子 G1、G3 和 G4 就已是共同不可满足的了。如果一些句子已经是共同不可满足的，那么再往其中添加额外的句子也无法使它们变成共同可满足的！）

有一点值得指出：你可能认为，一个论证只有在其前提是\emph{共同可满足}的时候才“有意义”。然而，我们对“什么是论证”以及“它何时有效”的定义，都不要求这一点。实际上，根据我们的定义，任何前提\emph{根本不可能同时为真}的论证，都是自动有效的！（练习：说服自己这是真的。）

\ifHTMLtarget
\section{必然真理、必然假与偶然性}
\else
\section[必然真理与必然假]{必然真理、必然假与偶然性}
\fi

在评估论证的有效性时，我们关心的是\emph{如果}前提为真，那么什么才会是真的。但有些句子本身\emph{就必定}为真。考虑这几个句子：


\begin{enumerate}
		\item\label{Acontingent} 天正在下雨。
		\item\label{Atautology} 天或者正在这里下着雨，或者没有。
		\item\label{Acontradiction} 天既正在这里下着雨，又没有在这里下着雨。
	\end{enumerate}


要知道 \cref*{Acontingent} 是否为真，你需要看看窗外或者查一下天气预报。它可能为真，也可能为假。一个能在不同情形下为真或为假的句子，被称为\define{偶然的（contingent）}。

\newglossaryentry{contingent sentence}
{
name=偶然的句子,
description={一个句子既不是\gls{necessary truth}也不是\gls{necessary falsehood}；它在某些情形下为真，在另一些情形下为假}
}

\Cref*{Atautology} 则不同。你不需要看窗外就知道它为真。无论天气怎样，天要么正在这里下雨，要么没有。这就是一个\define{必然真理（necessary truth）}。

\newglossaryentry{necessary truth}
{
name={必然真理},
description={一个在所有情形下都为真的句子}
}

同样，你也不需要查看天气来确定 \cref*{Acontradiction} 是否是真的。仅凭逻辑自身，它就必定为假。天可能正在这里下雨，而在镇子那边没下；可能现在正在下雨，但当你读完这句话时雨就停了；但天不可能在同一个地方、同一个时间既在下雨又没有下雨。所以，无论世界实际如何，天都不可能既在这里下雨又没有在这里下雨。它是一个\define{必然假（necessary falsehood）}。

\newglossaryentry{necessary falsehood}
{
name={必然假},
description={一个在所有情形下都为假的句子}
}

某个东西可能\emph{总是}为真，却仍然是偶然的。例如，如果宇宙从来不曾包含少于七个事物，那么“至少存在七个事物”这个句子就会一直为真。然而这个句子是偶然的：这个世界原本可能比现在小得多得多，那样的话这个句子就会为假。

\section{必然等值}

我们也可以考察两个句子\emph{之间}的逻辑关系。例如：


\begin{compactlist}
\item[] John 洗完盘子后去了商店。
\item[] John 去商店之前洗了盘子。
\end{compactlist}


这两个句子都是偶然的，因为 John 可能根本没去商店或洗盘子。然而，它们的真值必须相同。如果这两个句子中有一个为真，则两者都为真；如果有一个为假，则两者都为假。当两个句子在每一种情形下都具有相同的真值时，我们说它们是\define{必然等值（necessarily equivalent）}的。

\newglossaryentry{necessary equivalence}
{
name={必然等值},
text={必然等值},
description={一对句子具有的性质：在所有情形下，它们要么同为真，要么同为假}
}

\section*{逻辑概念总结}

\factoidbox{


\begin{factoiditems}
\item 一个论证是\define{有效的}，如果不存在前提全真而结论不为真的情形；否则它就是\define{无效的}。

\item 一个\define{必然真理}是在所有情形下都为真的句子。

\item 一个\define{必然假}是在所有情形下都为假的句子。

\item 一个\define{偶然的句子}既不是必然真理也不是必然假；它是一个在有些情形下为真，在另一些情形下为假的句子。

\item 两个句子是\define{必然等值}的，如果它们在所有情形下都同为真或同为假。

\item 一组句子是\define{共同可满足的}，如果存在一个情形使它们同时为真；否则就是\define{共同不可满足的}。
\end{factoiditems}


}

\practiceproblems
\problempart
\label{pr.EnglishTautology2}
对于下列各项，判断它是必然真理、必然假，还是偶然的？

\begin{compactlist}
\item Caesar 渡过了卢比孔河。
\item 有人曾渡过卢比孔河。
\item 没有人曾渡过卢比孔河。
\item 如果 Caesar 渡过了卢比孔河，那么有人曾渡过。
\item 尽管 Caesar 渡过了卢比孔河，但没有人曾渡过卢比孔河。
\item 如果任何人曾渡过卢比孔河，那么此人就是 Caesar。
\end{compactlist}

\problempart
对于以下各项，判断它是必然真理、必然谬误，还是偶然的？


\begin{compactlist}
\item 大象会在水中溶解。
\item 木材是一种轻便耐用的材料，可用于建造物品。
\item 如果木材是好的建筑材料，那么它可用于建造物品。
\item 我住在一栋三层高但只有两层楼高的建筑里。
\item 如果沙鼠是哺乳动物，那么它们会哺乳幼崽。
\end{compactlist}

\problempart
以下哪些句子对是必然等价的？

\begin{compactlist}
\item 大象会在水中溶解。 \\
	如果你把一头大象放进水里，它会分解。
\item 所有哺乳动物都会在水中溶解。 \\
	如果你把一头大象放进水里，它会分解。
\item George Bush 是第 43 任总统。 \\
	Barack Obama 是第 44 任总统。
\item Barack Obama 是第 44 任总统。 \\
	Barack Obama 是紧接着第 43 任总统之后的总统。
\item 大象会在水中溶解。 \\
	所有哺乳动物都会在水中溶解。
\end{compactlist}


\problempart
以下哪些句子对是必然等价的？

\begin{compactlist}
\item Thelonious Monk 演奏钢琴。 \\
	John Coltrane 演奏次中音萨克斯。
\item Thelonious Monk 与 John Coltrane 一起演出。 \\
	John Coltrane 与 Thelonious Monk 一起演出。
\item 所有职业钢琴演奏者都有大手。 \\
	钢琴演奏者 Bud Powell 有大手。
\item Bud Powell 患有严重的精神疾病。 \\
	所有钢琴演奏者都患有严重的精神疾病。
\item John Coltrane 有虔诚的宗教信仰。 \\
	John Coltrane 将音乐视为灵性的一种表达。
\end{compactlist}

\noindent \problempart
考虑以下句子：


\begin{compactlist}%[label=(\alph*)]
\item[G1.] \label{itm:at_least_four}野生动物园里至少有四只长颈鹿。
\item[G2.] \label{itm:exactly_seven} 野生动物园里恰好有七只大猩猩。
\item[G3.] \label{itm:not_more_than_two} 野生动物园里的火星人不超过两个。
\item[G4.] \label{itm:martians} 野生动物园里的每只长颈鹿都是火星人。
\end{compactlist}

现在考虑以下每组句子。哪些是共同可能的？哪些是共同不可能的？


\begin{compactlist}
\item 句子 G2、G3 和 G4
\item 句子 G1、G3 和 G4
\item 句子 G1、G2 和 G4
\item 句子 G1、G2 和 G3
\end{compactlist}

\problempart
考虑以下句子。

\begin{compactlist}%[label=(\alph*)]
\item[M1.] \label{itm:allmortal} 所有人都会死。
\item[M2.] \label{itm:socperson} Socrates 是人。
\item[M3.] \label{itm:socnotdie} Socrates 永远不会死。
\item[M4.] \label{itm:socmortal} Socrates 是会死的。
\end{compactlist}

哪些句子组合是共同可能的？将每项标记为“可能的”或“不可能的”。

\begin{compactlist}
\item 句子 M1、M2 和 M3
\item 句子 M2、M3 和 M4
\item 句子 M2 和 M3
\item 句子 M1 和 M4
\item 句子 M1、M2、M3 和 M4
\end{compactlist}

\problempart
\label{pr.EnglishCombinations2}
以下哪些情况是可能的？对于每种情况，如果可能，请给出一个例子；如果不可能，请解释原因。

\begin{compactlist}
\item 一个前提为一假一真的有效论证
\item 一个结论为假的有效论证
\item 一个有效论证，其结论是一个必然谬误
\item 一个无效论证，其结论是一个必然真理
\item 一个偶然的必然真理
\item 两个必然等价的句子，它们都是必然真理
\item 两个必然等价的句子，其中一个是必然真理，另一个是偶然的
\item 两个必然等价的句子，但它们共同是不可能的
\item 一个共同可能的句子集合，其中包含一个必然谬误
\item 一个共同不可能的句子集合，其中包含一个必然真理
\end{compactlist}

\problempart
以下哪些情况是可能的？对于每种情况，如果可能，请给出一个例子；如果不可能，请解释原因。

\begin{compactlist}
\item 一个有效的论证，其所有前提都是必然真理，而其结论是偶然的
\item 一个具有真前提和假结论的有效论证
\item 一个共同可能的句子集合，其中包含两个并非必然等价的句子
\item 一个共同可能的句子集合，其中的所有句子都是偶然的
\item 一个假的必然真理
\item 一个具有假前提的有效论证
\item 一对必然等价但并不共同可能的句子
\item 一个同时也是必然谬误的必然真理
\item 一个所有句子都是必然谬误的共同可能的句子集合
\end{compactlist}
